% CLAIM: the systems claims are about an algorithm that demonstrably trains.
% Evidence: experiments/countdown/results/*/README.md, plot_e13.py (F7),
% and the C6a-d sub-claims.
\section{Accuracy and task validation}
\label{sec:endtoend}

% ORDER (revised 2026-08-23, fifth review): determinism, then the synthetic
% estimator result, then the task prediction and its outcome, then the GRPO
% baseline, then a short pointer to the appendix stress test. The argument is
% prediction-before-campaign, so the prediction's evidence comes before the
% campaign's; stability moved to the appendix as exploratory.
This section is the validity chain the introduction promised, centered on a
real model (Qwen2.5-0.5B on the Countdown task) with a true-gradient
replication at 1.5B, all through the same instrument, in
an order where each claim rests on the one before it. First,
reproducibility: the systems prices of Section~\ref{sec:results} are
worth acting on only if what was measured is the algorithm, so the
determinism boundary is stated exactly. Second, quality: the estimator
study's prediction, then the campaign that tests it, closing the
unmeasured half of the low-rank belief. An exploratory stress test of
the stability folklore, a separate question, closes the section by
pointer; it lives in Appendix~\ref{sec:stability}.

\paragraph{Determinism and invariance on real hardware.} Two independent
20-generation runs of the same $D{=}8$ program agree byte for byte in every
logged reward and exactly in all 494M final parameters, and with the
fitness vector held fixed the $D{=}1$ and $D{=}8$ updates agree to
$6.3\times10^{-6}$ norm-relative error. The guarantee stops at separately
compiled programs, where XLA fusion and autotuning can flip near-tied
greedy tokens and the sampled trajectories then diverge; we document that
boundary rather than promise across it.

% F5 -- experiments/phase0/figures/f5-estimator-quality.png
% (experiments/phase0/plot.py --sigma 0.001 --shaping centered_ranks)
\begin{figure*}[t]
  \centering
  \includegraphics[width=\textwidth]{f5-estimator-quality.png}
  \caption{Estimator quality on the synthetic block, at the slice the task
    prediction uses ($\sigma{=}10^{-3}$, centered-rank shaping, the shaping
    the update applies): cosine between the ES update direction and the
    true gradient, against the population-to-effective-dimension ratio
    $N/d_{\mathrm{eff}}$, one panel per perturbation rank. Every fitted
    curve has slope one half. The design-based sampling schemes (orthogonal
    Hadamard, scrambled Sobol) are indistinguishable from plain mirrored
    sampling, the negative result Section~\ref{sec:related} records. The
    unpaired i.i.d.\ scheme sits a factor $1.34$--$1.38$ above the mirrored
    family at every matched population in the low-rank panels, the ordering
    Section~\ref{sec:results}'s baseline comparison cites. The low-rank
    panels pay a per-sample projection loss (a lower intercept) that their
    smaller effective dimension wins back at matched population, the
    cancellation the accuracy chain relies on.}
  \label{fig:f5}
\end{figure*}

\paragraph{But is the fast strategy accurate?} The speed results of
Section~\ref{sec:results} would mean little if the low-rank rewrite bought
its speed with estimation quality, so the accuracy chain is stated in
full, with its methods. The estimator study measured, on the
synthetic block of Section~\ref{sec:setup}, the cosine between the ES
update direction and the exact gradient from automatic differentiation,
30 replicates per cell over 456 cells, and fitted
$\log\cos = a + b\log(N/d_{\mathrm{eff}})$ per strategy; every fitted
slope is one half (Figure~\ref{fig:f5}).
Per sample, on those objectives, mirrored low-rank matches mirrored
full-rank almost exactly at matched population: each low-rank sample loses
about $2.2\times$ to projecting the gradient through thin factors, and
low-rank's smaller effective dimension wins almost exactly that factor back
(the cancellation behind Figure~\ref{fig:f5}). Those fitted curves,
evaluated at the task's operating point before the campaign ran,
predicted no rank ordering a three-seed campaign could resolve
(predicted cosines $1.7\times10^{-4}$ full against
$2.7$--$2.8\times10^{-4}$ for ranks 1 and 4, the low-rank family
slightly ahead). The
extrapolation is far and is stated as one: one to three decades below the
smallest measured $N/d_{\mathrm{eff}}$, from a synthetic block to a real
model, with rank 16 riding on the rank-1/rank-4 trend because the
estimator study never measured it. That extrapolation is now measured
rather than trusted. On Qwen2.5-0.5B itself, the cosine between one
full update and the true gradient of the same teacher-forced loss
(five seeds per cell, Table~\ref{tab:tb4}) lands within 10\% of the
fitted prediction for full rank, and uniformly at about half of it for
the low-rank family (ratios $0.43$--$0.67\times$ across ranks and
populations): the real transformer's loss surface charges roughly twice
what the synthetic block did for the low-rank projection, as one
systematic factor. Because the shortfall is uniform, the main relative
regularities transfer to the real model: the cross-rank tie at $N{=}30$, the
square-root scaling in $N$ inside every arm, and the unpaired scheme's
per-sample advantage, predicted at $1.34\times$ and measured at
$1.36\times$. Two further probes pin the correction down. Rerunning the
two bracketing arms on a freshly drawn frozen prompt batch reproduces
both ratios to two decimals ($0.94\times$ and $0.52\times$ on either
batch, Table~\ref{tab:tb6}), so neither shortfall is specific
to the original prompt batch. And on Qwen2.5-1.5B, with both predictions
frozen before the run, the scaling law transfers unrefitted (within-arm
$N$-scaling $2.80\times$ and $2.94\times$ against $\sqrt{8}=2.83$),
full rank stays on its raw prediction ($0.92$--$0.95\times$) at
$3.1\times$ the parameters, and the 0.5B-calibrated low-rank prediction
lands within $23$--$26\%$ where the raw fit misses by $2.4$--$2.6\times$;
the effective correction drifts from $0.53$ to about $0.40$ across the
scale step, so one constant carries most of the transfer without being
scale-invariant (Table~\ref{tab:tb6}). The mechanism's reading: in the
models, objective, population range, and perturbation ranks measured
here, the population is minuscule against the parameter count, and rank
was not the dominant limitation on estimation quality. That, arguably,
is \emph{why} the low-rank bet can work at all in this regime.

% generated by experiments/countdown/analysis_e15.py --latex; do not edit
\begin{table*}[t]
\centering
\small
\caption{The synthetic-to-real bridge, measured. Cosine between one
  production update and the true NLL gradient on Qwen2.5-0.5B
  ($\times 10^{-4}$; median [min, max] over 5 seeds), against the
  estimator study's fitted prediction at the same $N/d_{\mathrm{eff}}$
  (centered-rank slice, the shaping the update applies). The estimator
  study has no rank-16 curve, so that arm is measured only.}
\label{tab:tb4}
\begin{tabular}{llllr}
\toprule
arm & $N$ & measured & predicted & ratio \\
\midrule
full rank & 30 & $1.4$ [1.3, 1.9] & $1.6$ & $0.91\times$ \\
full rank & 240 & $4.2$ [3.8, 4.4] & $4.5$ & $0.94\times$ \\
rank 1 & 30 & $1.4$ [1.2, 1.7] & $2.7$ & $0.51\times$ \\
rank 1 & 240 & $3.9$ [3.8, 4.2] & $7.6$ & $0.52\times$ \\
rank 4 & 30 & $1.9$ [0.7, 2.0] & $2.8$ & $0.67\times$ \\
rank 4 & 240 & $4.4$ [4.1, 4.8] & $7.9$ & $0.56\times$ \\
rank 16 & 30 & $1.5$ [1.1, 1.7] & -- & -- \\
rank 16 & 240 & $4.1$ [3.9, 4.5] & -- & -- \\
rank 1, unpaired & 30 & $1.6$ [1.1, 1.9] & $3.8$ & $0.43\times$ \\
rank 1, unpaired & 240 & $5.3$ [5.1, 5.5] & $10.6$ & $0.50\times$ \\
\bottomrule
\end{tabular}
\end{table*}

% generated by experiments/countdown/tb6_e16.py --latex; do not edit
\begin{table}[t]
\centering
\small
\setlength{\tabcolsep}{3.2pt}
\caption{Cosine between one ES update and the exact gradient of the
  next-token loss ($\times 10^{-4}$, median over 5 perturbation seeds), at equal
  population; \emph{1 unp.} is rank 1 without mirrored pairs.
  \emph{ref.}: the mirrored full-rank fit from the synthetic block
  (Figure~\ref{fig:f5}) at the same $N/P$, a post hoc reference; unpaired
  sampling should sit about $\sqrt2$ above it. Batch 2 repeats two
  configurations on freshly drawn prompts. Ranges and the preregistered
  predictions are in Table~\ref{tab:tb6-audit}.}
\label{tab:tb6}
\begin{tabular}{lrrrrrrr}
\toprule
 & & \multicolumn{5}{c}{rank} & \\
\cmidrule(lr){3-7}
model & $N$ & full & 1 & 4 & 16 & 1 unp. & ref. \\
\midrule
0.5B & 30 & $1.42$ & $1.37$ & $1.87$ & $1.49$ & $1.60$ & $1.56$ \\
0.5B & 240 & $4.23$ & $3.91$ & $4.42$ & $4.09$ & $5.33$ & $4.48$ \\
0.5B, batch 2 & 240 & $4.20$ & $3.91$ & -- & -- & -- & $4.48$ \\
1.5B & 30 & $0.83$ & $0.77$ & -- & -- & -- & $0.87$ \\
1.5B & 240 & $2.32$ & $2.27$ & -- & -- & -- & $2.51$ \\
\bottomrule
\end{tabular}
\end{table}


\paragraph{The task, and what was fixed in advance.} The campaign that
tests the prediction is specified in Table~\ref{tab:e2e}: both algorithm
families train on the same puzzle pool with the published core
hyperparameters of \citet{qiu2025} (one disclosed budget difference:
8 puzzles per update against their protocol's 200 training samples), are
scored on the same 2000 disjoint held-out puzzles by greedy decode, and
stop at the same total budget of scored training completions, which is
the unit the x-axis of Figure~\ref{fig:f7} counts.

\begin{table*}[t]
\centering
\small
\caption{End-to-end setup. A sample evaluation is one scored
  \emph{training} completion (held-out scoring is not counted); both
  families consume 240 per update (ES: 30 members $\times$ 8 puzzles;
  GRPO: 30 completions $\times$ 8 prompts), so the 120{,}000 budget is
  500 updates for each.}
\label{tab:e2e}
\begin{tabular}{lp{0.78\linewidth}}
\toprule
task, reward & Countdown; 1.0 solved, 0.1 well-formed but wrong, 0.0
  otherwise; ``solved'' means reward 1.0 \\
data & training pool 256 puzzles, 8 per update; held-out 2000 disjoint
  puzzles \\
prompts, decode & chat template, prompts padded to 128 tokens, 96 new
  tokens; ES rollouts and all held-out evals decode greedily \\
ES arms & $N{=}30$, $\sigma{=}10^{-3}$, $\alpha{=}5\times10^{-4}$ (the
  published values; implemented as $\mathrm{lr}{=}\alpha\sigma{=}5\times10^{-7}$
  because the library's update divides by $\sigma$), centered-rank
  shaping; full rank and ranks 1, 4, 16, plus rank 1 with the embedding
  (27\% of parameters) frozen \\
GRPO arm & TRL implementation; group size 30, 8 prompts per step,
  learning rate $10^{-5}$ constant, KL $\beta{=}10^{-3}$, clip
  $\epsilon{=}0.2$, rollout temperature 1.0 with no top-$p$/top-$k$ \\
budget, seeds & 120{,}000 sample evaluations per run; seeds 0--2 per arm \\
evaluation & 2000 held-out puzzles every 50 updates, plus before the
  first and after the last \\
hardware & one A100-SXM4-80GB; the campaign is $\sim$13\,h \\
\bottomrule
\end{tabular}
\end{table*}

% F7 -- experiments/countdown/figures/f7-e13-heldout.png (plot_e13.py)
\begin{figure*}[t]
  \centering
  \includegraphics[width=0.72\textwidth]{f7-e13-heldout.png}
  \caption{Countdown (arithmetic puzzle solving; the task shape of
    TinyZero, \citealp{tinyzero2025}) with Qwen2.5-0.5B \citep{qwen25}:
    held-out reward against training sample evaluations, mean over three
    seeds,
    min--max band. The four ES perturbation ranks lie on top of each
    other; freezing the embedding shows no detectable difference in this
    three-seed experiment; the GRPO reference with its published settings
    is seed-unstable: two seeds plateau at 4.9\% of held-out puzzles
    solved and one never solves any, though none falls below the
    base-model floor. The two labeled base-model floors are the same
    weights under the JAX and HuggingFace decoders, apart by a bfloat16
    padding artifact, which bounds how small a cross-family gap can
    honestly be read from this figure.}
  \label{fig:f7}
\end{figure*}

\paragraph{Task quality across ranks.} At population $N{=}30$, held-out
performance shows no detectable difference across full-rank and
rank-\{1,4,16\} perturbations in this three-seed experiment
(Figure~\ref{fig:f7}); no inferential test is claimed at $n{=}3$. The
outcome is the one the estimator study's fitted curves predicted, which
put all ranks near parity at this
population-to-effective-dimension ratio, closer than a three-seed
campaign resolves; a non-difference is consistent
with the estimator mechanism, not proof of it.

% F7b -- experiments/countdown/figures/f7b-e13-wallclock.png (plot_e13.py)
\begin{figure*}[t]
  \centering
  \includegraphics[width=0.72\textwidth]{f7b-e13-wallclock.png}
  \caption{The curves of Figure~\ref{fig:f7}, ES arms only, against
    cumulative accelerator time in steady state: generations 0 and 1 are
    both excluded because both are compilation-scale events in every arm
    and seed (335--671\,s against steady-state 2.4--4.3\,s per update),
    and held-out evaluation decode is excluded on both axes as
    measurement overhead. The evaluation at generation $g$ runs before
    update $g$, so its x is the cumulative time of updates strictly
    before $g$. Each point's x is the mean cumulative time over three
    seeds; bands are min--max of reward. GRPO is deliberately absent: it
    runs in a different framework with a different decoder, so its wall
    clock is not commensurate with these curves.}
  \label{fig:f7b}
\end{figure*}

\paragraph{Equal samples, unequal seconds.} Figure~\ref{fig:f7}
resolves no rank-dependent quality difference over three seeds;
Figure~\ref{fig:f7b} prices the same curves in accelerator time. A
steady-state full-rank update costs 4.29\,s against 2.41\,s for rank 1,
means over the 498 steady-state updates each run observed (the
regeneration compute the cost surfaces of Section~\ref{sec:results}
price at microbenchmark scale, here on the real model): a $1.78\times$
per-update advantage, an estimated 2147\,s against 1205\,s for a
500-update run. Time to first reach the shared plateau (held-out reward
0.15) is 419--1069\,s for full rank against 236--483\,s for rank 1
across seeds; the ranges overlap between 419 and 483\,s, so three seeds
do not resolve a strict time-to-threshold ordering, and the per-update
ratio is the clean claim. This is where the systems and quality halves
meet: the cheap perturbation costs $1.78\times$ less per update while
Figure~\ref{fig:f7} shows it needed no more samples.

\paragraph{The reference baseline, honestly.} GRPO
\citep{deepseekmath2024} with the published hyperparameters of
\citet{qiu2025}, untuned in either direction: two of three seeds plateau
early at 4.9\% of held-out puzzles solved, and one never solves a
held-out puzzle, flat at format-only reward from step 50. None falls
below its decoder's base-model floor, so none meets the pre-specified
collapse definition of the appendix stress test; the word collapse is
reserved for runs that do. We report this as seed variance in the
reference, not as a
victory for ES; the figure's role is that both algorithm families run end
to end in comparable budgets, not a leaderboard.

\paragraph{An exploratory robustness stress test.} The folk belief that
ES does not need RL's stabilizers is a question about comparative
optimizer robustness, separate from this paper's systems argument, so
its full treatment lives in Appendix~\ref{sec:stability}. The scoped
summary: with metrics and falsification criteria fixed in a committed
protocol before the data, this GRPO recipe at $8\times$ its published
learning rate collapsed to zero reward on all three seeds while ES
collapsed at neither tested endpoint of the same 64-fold interval; our
own criteria caught the part of the belief that was wrong, since
removing the recipe's individual stabilizers at published settings
changed little; and the surviving ES runs drifted far from the
pretrained weights, a retention cost reported beside the robustness. The
multipliers are per-family and the interval interior was not explored,
so this is a stress test's evidence at one task, one horizon, and three
seeds, not a stability margin.
